\begin{table}[!htbp]
\centering
\manualtablecaption{Table S7. Diet composition of the scenario baskets and UPF targets used in the analyses}
\footnotesize
\resizebox{\textwidth}{!}{%
\begin{tabular}{lrrrrr}
\toprule
\multicolumn{6}{l}{\textbf{Panel A. Diet composition of replacement baskets (energy share, \%)}}\\
\textbf{Food group} & \textbf{\shortstack{Observed\\non-UPF}} & \textbf{NOVA1} & \textbf{NOVA3} & \textbf{\shortstack{NOVA1/NOVA3\\mix}} & \textbf{NDG target}\\
\midrule
Fruits & 7.3 & 13.7 & 5.7 & 9.0 & 18.0\\
Vegetables (outdoor) & 3.3 & 7.8 & 1.4 & 4.0 & 25.0\\
Legumes and pulses & 0.4 & 1.2 & 0.1 & 0.5 & 5.0\\
Cereals (without rice) & 8.0 & 7.1 & 2.0 & 4.0 & 6.0\\
Rice & 4.5 & 13.0 & 0.0 & 5.3 & 3.0\\
Pasta & 6.2 & 0.0 & 13.2 & 7.8 & 3.0\\
Baked products & 25.8 & 0.0 & 55.4 & 32.8 & 3.0\\
Starchy vegetables & 0.2 & 0.7 & 0.0 & 0.3 & 3.0\\
Dairy products & 10.0 & 5.0 & 15.4 & 11.1 & 13.0\\
Beef & 3.3 & 7.1 & 0.6 & 3.3 & 4.0\\
Lamb and mutton & 0.3 & 1.0 & 0.0 & 0.4 & 0.5\\
Other meats & 0.0 & 0.0 & 0.0 & 0.0 & 0.5\\
Poultry & 7.4 & 21.6 & 0.0 & 8.8 & 2.5\\
Pork & 6.1 & 10.5 & 5.2 & 7.4 & 1.0\\
Fish and seafood & 0.9 & 1.8 & 0.5 & 1.0 & 3.0\\
Eggs & 2.9 & 8.3 & 0.0 & 3.4 & 1.5\\
Oil crops & 13.4 & 1.2 & 0.4 & 0.7 & 6.0\\
Nuts and seeds & 0.1 & 0.0 & 0.2 & 0.1 & 2.0\\
\bottomrule
\end{tabular}
}

\vspace{0.5em}
\begin{tabular*}{\textwidth}{@{\extracolsep{\fill}}p{7.0cm}rr@{}}
\toprule
\multicolumn{3}{l}{\textbf{Panel B. Scenario-level UPF targets}}\\
\textbf{Scenario} & \textbf{\shortstack{Target UPF\\(\% energy)}} & \textbf{\shortstack{Reduction vs\\baseline (\%)}}\\
\midrule
Observed baseline & 16.66 & 0.0\\
10\% proportional reduction & 14.99 & 10.0\\
20\% proportional reduction & 13.33 & 20.0\\
50\% proportional reduction & 8.33 & 50.0\\
NDG conservative target & 5.45 & 67.3\\
NDG central target & 3.63 & 78.2\\
NDG strict target & 1.81 & 89.1\\
\bottomrule
\end{tabular*}
\normalsize
\tabnote{\textbf{Panel A.} Each column reports the share of total energy contributed by each food group \textit{within} that basket; columns sum to 100\%. \textbf{Observed non-UPF} is the actual energy composition of all non-UPF foods consumed (NOVA~1, 2, and 3 combined) by adults aged 30--69 in ENNyS~2, restricted to foods with environmental coefficients and excluding plain drinking water. This column is an empirical reference and is \textit{not} used as a replacement basket in the counterfactual scenarios. \textbf{NOVA~1} and \textbf{NOVA~3} show the energy composition within minimally processed and processed non-UPF foods, respectively, based on observed ENNyS~2 intakes. In the current Argentine data, the NOVA~1 basket is concentrated in animal-source foods together with fruits and rice, whereas the NOVA~3 basket is dominated by baked products, dairy, pasta, and other cereal-based items. \textbf{NOVA~1/NOVA~3 mix} is a weighted average of those two baskets, with weights derived from their respective model-mapped energy totals (approximately 39\% NOVA~1 and 61\% NOVA~3). \textbf{NDG target} shows the normative food-group proportions prescribed by Argentina's National Dietary Guidelines (NDG~2016)\textemdash these are guideline recommendations, not observed intakes, and differ substantially from the current diet (notably higher in fruits and vegetables, lower in baked products and oil crops). \textbf{Panel B.} ``Target UPF (\% energy)'' is the counterfactual share of ultra-processed foods in total dietary energy; ``Reduction vs baseline (\%)'' is the proportional reduction from the observed baseline reported in Panel~B. Proportional scenarios (10\%, 20\%, 50\%) scale the UPF energy share downward uniformly without altering the food-group composition of the non-UPF diet. NDG scenarios impose the NDG basket composition from Panel~A and derive the residual UPF allowance from the guideline food-group targets, yielding substantially larger reductions (67--89\%) because Argentina's current UPF share is far above the level implied by guideline-adherent dietary patterns. The three NDG profiles (conservative, central, strict) correspond to progressively tighter guideline adherence, yielding target UPF shares of 5.4\%, 3.6\%, and 1.8\%, respectively.}
\end{table}
